\section{A Formal Model of Governance Drift (RQ2, RQ4)}
\label{sec:model}

\subsection{States and Contexts}
\label{sec:model-states}

Fix a deployment $x$ into environment $e$ with approval events at
$t_0<\cdots<t_k\leq t$.

\textbf{Approved-snapshot history.} Per \cref{def:approved}, every approval
event appends an immutable snapshot
\[
\begin{aligned}
\Gapp^{x}(t)&=\{G_{\mathrm{app}}^{(0)}(t_0),\ldots,
G_{\mathrm{app}}^{(k)}(t_k)\},\\
G_{\mathrm{app}}^{(i)}&=\bigl(m_i,\delta_i,\dig(\pi_i),A_i^{\mathrm{proof}},
\iota_i,\sigma_i\bigr).
\end{aligned}
\]
Each record contains the approved manifest content $m_i$ and artifact digests $\delta_i$ it
resolved to at approval; the pinned policy version; the authorization basis
$A_i^{\mathrm{proof}}$ (the immutable proof that approval instruments covered
the named subjects and scope at their execution times); the intent linkage $\iota_i$ (change
record, revision); and the environment assumptions $\sigma_i$ (the declared
properties of $e$---identity scopes, exposure---under which risk was
assessed). Let $b(t)$ select the greatest $i$ such that $t_i\leq t$ and the
snapshot scope covers $(x,e)$. The admitted basis is
\[
B_x(t)=G_{\mathrm{app}}^{(b(t))}
=\bigl(m_B,\delta_B,\dig(\pi_B),A_B^{\mathrm{proof}},
\iota_B,\sigma_B\bigr).
\]
Selection does not require current validity: an expired or revoked basis must
remain selected so its divergence can be reported. Superseded snapshots
remain available for audit and lineage; benign reapproval appends a new basis
rather than mutating history. Append-only means records cannot be rewritten,
not that their storage cannot be lost; loss of $A_B^{\mathrm{proof}}$ is an
evidence failure, as S9 exercises.

\textbf{Observed state.} $\Gobs(t) = (m_t, \delta_t, \varepsilon_t)$: the
observed configuration, the digests actually executing, and the observed
environment properties needed to evaluate the predicates in $\sigma_B$.

\textbf{Contexts.} Three time-indexed context streams:
$\Pol(t)$, the policy version(s) governing $e$ at $t$ with their evaluation
semantics; $\Auth_{\mathrm{live}}(t)$, the authenticated current administrative
state of instruments (revocations, continuing validity, expiry); and $\Int(t)$, the intent records and the
desired-state revision currently driving delivery. Evidence availability,
where needed, is a fourth stream $\Evd(t)$ (the records from which the above
can be established).

\textbf{Authorization semantics.} Every instrument $a$ declares
$\mathrm{mode}(a)\in\{\text{one-shot},\text{continuing},
\text{temporary-exception}\}$, subject, scope, validity interval, and
revocation policy. One-shot authority is required at execution and its proof
is retained afterward; continuing authority must remain valid while its
effect persists; a temporary exception is continuing only until its expiry.
Revocation is prospective unless the record explicitly declares retroactive
invalidation. Let $\mathrm{NeedsLive}(a)$ hold for continuing and temporary
instruments, and for a one-shot instrument only when its declared revocation
policy permits retroactive invalidation. For an effect executed at $\tau$,
let $\mathrm{AppObs}(a,t)$ mean that its proof is available and, whenever
$\mathrm{NeedsLive}(a)$ holds, authenticated live status is also available.
Write $V_1$ for validity at $\tau$ without retroactive revocation, $V_c$
for current validity without revocation, and $V_e$ for $V_c$ while the
declared exception has not expired. Define the three-valued applicability function
\[
\mathrm{Applicable}(a,\tau,t)=
\begin{cases}
\bot,&\neg\mathrm{AppObs}(a,t);\\
V_1(a,\tau,t),&\text{one-shot};\\
V_c(a,t),&\text{continuing};\\
V_e(a,t),&\text{temporary exception}.
\end{cases}
\]
A retained historical proof is therefore not confused with currently
effective authority. Missing proof or live status yields undecidable where
the selected mode requires it; an authenticated revocation yields a polar
inconsistent verdict.

\textbf{Legitimate successors.} Let
$\mathrm{Succ}(\iota_B,t)$ be the set of revisions $\iota_n$ for which
there exists a finite chain $\iota_B,\iota_1,\ldots,\iota_n$ with transition
execution times $\tau_1,\ldots,\tau_n$, such that every edge
$(\iota_{j-1},\iota_j)$ and its artifact digest are covered by an instrument
$a_j$ for which $\mathrm{Applicable}(a_j,\tau_j,t)$ holds. A revision is a
legitimate successor iff it belongs to this set.
This definition excludes mere Git ancestry: lineage without a valid approval
edge does not extend the admitted basis.

\subsection{The Governance-Divergence Vector}
\label{sec:model-distance}

We define one divergence component per substantive relation; each is a function of exactly the
inputs its component needs, which is what makes the detectability results of
\cref{sec:study} possible to state.

\begin{itemize}[leftmargin=1.5em,itemsep=1.5pt,topsep=2pt]
\item $\dconf(t)$: a normalized difference over the managed configuration
graph between desired$(t)$ and $m_t$ (field-level graph diff with
runtime-owned fields factored out; \cref{sec:study} shows the normalization
is not cosmetic but the difference between signal and noise). This is an
operational consistency component included to give the evaluator a complete
control-state report; nonzero $\dconf$ alone does not imply that either the
desired or observed state lacks authorization coverage.
\item $\dpol(t)$: the set of controls in $\Pol(t)$ that the running state
violates although $B_x(t)$ records compliance under $\pi_B$---empty iff admission would
re-succeed under today's policy. Reported as a violated-control set, not a
number: policy versions form a revision structure, not a metric space, and
``how far'' is policy-weighted judgment, not geometry.
\item $\dauth(t)$: the set of relied-upon instruments for which
$\mathrm{Applicable}(a,\tau,t)$ is false, plus running digests in $\delta_t$
not covered by an applicable instrument grounded in $A_B^{\mathrm{proof}}$
and $\Auth_{\mathrm{live}}(t)$. If applicability is $\bot$, the authorization
mask entry is zero and the substantive component is undecidable rather than
included in this set.
\item $\dint(t)$: the mismatch between the revision driving delivery and the
revision the valid authorization covers: the current revision is outside
$\{\iota_B\}\cup\mathrm{Succ}(\iota_B,t)$.
\item $\denv(t)=\{q\in\sigma_B:q(\varepsilon_t)=\mathrm{false}\}$. The laboratory's
declared IAM-scope and exposure fields instantiate equality predicates, but
the definition permits ranges, membership, and monotone safety predicates.
\end{itemize}

Evidence availability is ontologically different from those distances. Define
the observability mask
$\ObsMask(t)=(o_{\mathrm{conf}},o_{\mathrm{pol}},o_{\mathrm{auth}},
o_{\mathrm{int}},o_{\mathrm{env}})$, where $o_j=1$ iff every required item in
the component's evidence contract is present, complete, fresh within its
declared bound, integrity-valid, source-authenticated, linked to subject
$x$, and delivered by a healthy stream. Confidence is therefore not an
unscoped probability: it is a named contract and threshold published with
the evaluator. A missing or contract-failing record yields \emph{undecidable};
an authenticated record stating ``revoked,'' ``negative,'' or ``never
issued'' is available evidence and may yield \emph{inconsistent}. Evidence drift is a
transition of any required $o_j$ from one to zero. Each substantive component
therefore has three-valued reporting semantics---consistent, inconsistent, or
undecidable---but only the first two are values of the substantive relation.

\begin{definition}[Governance Divergence Vector; Consistency]
\label{def:gd}
$\GDsub(t) = \bigl(\dconf, \dpol, \dauth, \dint, \denv\bigr)(t)$ and
$\GD(t)=(\GDsub(t),\ObsMask(t))$. Operational consistency
$\mathsf{OpCons}(x,t)$ holds iff $\dconf=0$ and its evidence is decidable.
Governance coverage $\mathsf{GovCov}(x,t)$ holds iff
$\dpol=\dauth=\dint=\denv=0$ and their mask entries are one. Total
$\Cons(x,t)=\mathsf{OpCons}(x,t)\land\mathsf{GovCov}(x,t)$. Onset of governance drift is
$\inf\{t > t_0 : \lnot\Cons(x,t)\}$; \emph{dwell time} is the duration of
inconsistency; \emph{detection latency} is detection time minus onset.
\end{definition}

\begin{figure}[t]
\centering
\resizebox{\columnwidth}{!}{% Drift lifecycle timeline (single column, resized)
\begin{tikzpicture}[x=1mm, y=1mm]
  \def\W{92}
  \draw[flow] (0,0) -- (\W,0) node[lbl, right]{$t$};
  \foreach \x/\lab/\col in {6/{approval $t_0$}/black!70,
                            30/{drift onset}/red!70,
                            56/{detection}/orange!90!black,
                            74/{remediation}/green!55!black}{
    \draw[\col] (\x,0) -- (\x,2.5);
    \node[lbl, above, text=\col] at (\x,3) {\lab};
  }
  \draw[decorate,decoration={brace,mirror,amplitude=3pt}, black!60]
    (30,-2) -- (56,-2) node[midway, below=3pt, lbl] {detection latency (DDL)};
  \draw[decorate,decoration={brace,mirror,amplitude=3pt}, black!60]
    (30,-8.5) -- (74,-8.5) node[midway, below=3pt, lbl] {dwell time (exposure)};
  \node[lbl, below=1mm] at (18,-13.5) {\smlbl consistent};
  \draw[line width=2.2pt, green!55!black] (6,-13) -- (30,-13);
  \draw[line width=2.2pt, red!75] (30,-13) -- (74,-13);
  \draw[line width=2.2pt, green!55!black] (74,-13) -- (\W-4,-13);
  \node[lbl, below=1mm, text=red!75] at (52,-13.5) {\smlbl inconsistent (classed by $\GD$)};
  \node[lbl, below=1mm] at (84,-13.5) {\smlbl restored};
  \node[lbl, text=black!60] at (46,8)
    {onset is class-specific: an event ($\pi$ revision, re-tag), a clock (expiry), or a convergence (rollback)};
\end{tikzpicture}
}
\caption{The drift lifecycle and its temporal quantities
(\cref{def:gd}). Onset is class-specific---an organizational event (policy
revision), a clock crossing (exception expiry), an exogenous event
(registry re-tag), or a reconciler convergence (rollback). Detection latency
is the tier-dependent quantity measured in \cref{sec:study}; dwell time is
the exposure window that remediation closes, with the remediation family
determined by the drifted component (\cref{sec:model-vector}).}
\label{fig:lifecycle}
\end{figure}

\subsection{Why the Distance Is a Vector, Not a Scalar}
\label{sec:model-vector}

A scalar $\GD$ was the natural first proposal, and we reject it as primitive
for three reasons. \emph{Incommensurability}: the components take values in
different structures (graph diffs, control sets, validity states,
three-valued verdicts); any scalarization imports policy weights, and those
weights are governance choices, not properties of the system.
\emph{Remediation asymmetry}: the components have different remediations---%
reconcile ($\dconf$), re-evaluate and re-approve or migrate ($\dpol$),
re-authorize or terminate the exception ($\dauth$), restore or re-approve
intent ($\dint$), restore assumptions or re-assess ($\denv$), repair evidence
where $\ObsMask$ is zero---so collapsing them destroys exactly the information an operator
needs. \emph{Temporal structure}: components drift on different clocks
(policy revisions are discrete organizational events; exception expiry is a
scheduled instant; registry re-tags are exogenous), so a scalar time series
conflates processes that must be attributed separately. A weighted composite
for reporting is legitimate \emph{on top of} the vector (\cref{sec:metrics}),
with published weights, exactly as a risk score summarizes but does not
replace findings.

\subsection{Detectability, Formally}
\label{sec:model-detect}

Each component divergence is a function of specific streams:
$\dconf$ of (desired, $\Gobs$); $\dpol$ additionally of $\Pol(t)$ and
$\pi_B$; $\dauth$ of $A_B^{\mathrm{proof}}$, $\Auth_{\mathrm{live}}(t)$,
execution times, and $\delta_t$; $\dint$ of $\Int(t)$ and $\iota_B$;
$\denv$ of $\varepsilon_t$ and $\sigma_B$; and every
entry of $\ObsMask$ of the evidence needed for its corresponding computation.
\Cref{prop:impossible} is the negative direction of this
dependency structure; the positive direction is equally direct:

\begin{proposition}[Tier sufficiency]
\label{prop:tiers}
A detector with access to the streams a component depends on decides that
component's consistency (in the closed-world model, exactly; in the field, up
to the fidelity of the streams). Consequently a sufficient detector for full
governance consistency joins the configuration comparison with the policy,
authorization, intent/lineage, and environment streams---the tiered
architecture of \cref{sec:architecture}, whose scenario-to-tier realization
\cref{sec:study} checks cell by cell.
\end{proposition}
\begin{proof}
Each component divergence in \cref{sec:model-distance} is defined as a
computable function of the named streams; providing the streams provides the
function's inputs.
\end{proof}

The proposition is an input-enumeration result; whether a concrete
implementation realizes it is a separate, checkable question---%
\cref{sec:study} performs exactly that check on our released reference
semantics, and \emph{only} that check: agreement there validates the
implementation, not the model.

\begin{proposition}[Tier minimality]
\label{prop:tier-minimality}
For every stream added at T1--T4, there exist two worlds that are
indistinguishable to every lower tier but have different governance verdicts
for a component first decidable from that stream. Therefore no lower tier
decides that component in general.
\end{proposition}
\begin{proof}
Construct the pairs while holding all lower-tier observations fixed. For T1,
keep desired and observed configuration identical and let the current policy
admit the deployment in one world but reject it in the other. For T2, keep the
T1 projection fixed and let the same effective exception be valid in one world
but expired in the other; alternatively, retain versus delete its approval
record, changing the authorization/intent mask from decidable to undecidable.
For T3, keep the manifest tag and every T2 stream fixed but let the running tag
resolve to an approved digest in one world and an uncovered digest in the
other. For T4, keep all T3 streams fixed and change only an unmanaged
environment property recorded in $\sigma_B$. Each pair has identical
lower-tier projections and different component verdicts, so any lower-tier
detector returns the same value in both and errs in one.
\end{proof}

Unlike \cref{prop:tiers}, this is a necessity result: the staircase is not
merely a convenient implementation decomposition. Its scope remains relative
to the declared component semantics; adding a genuinely equivalent stream at
a lower tier widens that tier's projection rather than refuting the result.

Two modeling notes bound the claims. First, \emph{granularity}: whether the
registry re-tag scenario is ``configuration'' drift depends on reference
semantics---under tag semantics the configuration is unchanged; under digest
semantics it changed. The model takes no side: it requires only that
$B_x(t)$ record resolved digests, which relocates the question from metaphysics
to evidence (what did the approval bind?), where \cref{sec:study} shows it is
decidable at the lineage tier. Second, \emph{legitimate context change}:
policy supersession does not make drift \emph{culpable}---the deployment was
legitimate under $\pi_7$ and the organization changed the rules. The model
deliberately measures divergence, not blame; whether policy drift triggers
migration, exemption, or re-approval is a governance decision the vector
informs (\cref{sec:discussion}).
